% =============================================================================
%  Seccion: Diccionario de datos
%
%  Version en LaTeX de docs/Diccionario_Sima.md. Si cambia el consolidado,
%  hay que actualizar AMBOS archivos (el .md es el que se consulta al programar).
%
%  Este archivo NO lleva preambulo: main.tex lo mete con \input. Solo usa
%  paquetes que main.tex ya carga (booktabs, amsmath, cleveref). Si hace falta
%  uno nuevo, se pide en una issue de main.tex; no se agrega aqui.
%
%  Las tablas anchas usan el entorno table* (flotante a dos columnas): en un
%  documento 5p una tabla de 6 o 7 columnas no cabe en una sola columna. Los
%  flotantes a dos columnas SOLO pueden ir arriba de la pagina, por eso llevan
%  [t] y no [h].
% =============================================================================

\section{Diccionario de datos}\label{sec:diccionario}

Esta seccion documenta el consolidado que produce
\texttt{src/importar\_datos.py} a partir de los seis libros anuales del SIMA
(\texttt{data/raw/BD 2020.xlsx} a \texttt{data/raw/BD 2025.xlsx}): que contiene
cada variable, como estan cubiertas las estaciones y que anomalias de calidad
hay que atender antes de cualquier analisis multivariado.

\subsection{Dimension del conjunto de datos}\label{sec:dic:dimension}

El consolidado tiene 754\,603 filas y 18 columnas: 3 identificadoras y 15
parametros medidos. Cada fila es una estacion en una hora, la cobertura va del
1 de enero de 2020 a las 00:00 al 31 de diciembre de 2025 a las 23:00, y no hay
registros duplicados. De las 11\,319\,045 celdas de parametros, 1\,188\,343
estan vacias: un 10.50\,\% de faltantes.

En el origen cada libro anual trae una hoja por estacion --- 13 en 2020, 14 en
2021 y 15 de 2022 en adelante --- en formato ancho: una fila por hora y una
columna por parametro. No existe una columna de estacion; la consolidacion
promueve el nombre de la hoja a la columna \texttt{estacion} y apila los seis
anios.

\subsection{Variables identificadoras}\label{sec:dic:identificadoras}

\begin{description}
  \item[\texttt{fecha}] Marca temporal de la observacion, con resolucion
    horaria (\texttt{datetime64}), de 2020-01-01 00:00 a 2025-12-31 23:00. Sin
    faltantes.
  \item[\texttt{estacion}] Categorica nominal de 15 niveles: CE, NE, NE2, NE3,
    NO, NO2, NO3, NTE, NTE2, SE, SE2, SE3, SO, SO2 y SUR. Sin faltantes.
  \item[\texttt{anio}] Anio calendario, derivado del archivo de origen.
    Numerica discreta (o categorica ordinal), con valores de 2020 a 2025. Sin
    faltantes.
\end{description}

\subsection{Contaminantes criterio}\label{sec:dic:contaminantes}

Los ocho contaminantes criterio son numericos continuos y se resumen en
\cref{tab:dic-contaminantes}. El rango de operacion que se reporta abarca desde
el anio mas restrictivo hasta el mas permisivo del periodo, porque el limite
cambia anio con anio (\cref{sec:dic:centinela}). Segun \texttt{Etiquetas.xlsx},
O$_3$, SO$_2$ y NO$_2$ pueden pasarse a ppm dividiendo entre 1000.

\subsection{Parametros meteorologicos}\label{sec:dic:meteorologicos}

Los siete parametros meteorologicos aparecen en \cref{tab:dic-meteo}. Dos de
ellos no pueden tratarse como variables continuas ordinarias:

\begin{itemize}
  \item \texttt{WDR} es una variable \textbf{circular}, aunque se almacene como
    numero: $0^{\circ}$ y $360^{\circ}$ son la misma direccion (norte). No
    puede promediarse, escalarse ni correlacionarse como las demas; hay que
    descomponerla en $\sin\theta$ y $\cos\theta$, o bien en las componentes de
    viento $u = -\,\mathrm{WSR}\cdot\sin\theta$ y
    $v = -\,\mathrm{WSR}\cdot\cos\theta$. Es un atributo derivado que debe
    construirse en la etapa de transformacion.
  \item \texttt{RAINF} y \texttt{SR} son de \textbf{cero inflado}: el 98.4\,\%
    de los registros validos de precipitacion y el 39.6\,\% de los de radiacion
    solar valen exactamente cero (noche y dias sin lluvia). Son ceros
    legitimos, no faltantes ni errores, pero rompen los supuestos de normalidad
    y distorsionan cualquier estandarizacion.
\end{itemize}

\subsection{Catalogo de estaciones}\label{sec:dic:estaciones}

\Cref{tab:dic-estaciones} lista las 15 estaciones y su cobertura calculada
sobre el consolidado. Dos estaciones no estan documentadas en el catalogo del
SIMA y ademas entraron tarde a la red: NE3 aparece a partir de 2021 y NO3
hasta diciembre de 2022, de modo que cualquier analisis que exija el panel
completo debe excluirlas o recortarse al periodo comun.

\subsection{Faltantes}\label{sec:dic:faltantes}

\Cref{tab:dic-faltantes} da el porcentaje de nulos por variable, que va del
3.66\,\% de \texttt{SR} al 25.23\,\% de \texttt{PM2.5}. Ese promedio esconde lo
importante: el faltante no esta repartido de forma pareja entre estaciones ni
entre anios. Para el ozono, 2020 es practicamente inservible: varias
estaciones (NTE, SUR, SE2, NE, NO) superan el 98\,\% de faltante ese anio; de
2022 en adelante el faltante de O$_3$ baja a entre 1\,\% y 8\,\% en casi toda
la red.

Los faltantes de este conjunto \textbf{no son aleatorios} (no son MCAR): el
SIMA invalida horas por calibracion, falla o condicion anomala antes de
entregar los archivos, y esas horas quedan como celda vacia sin que el motivo
sea recuperable. Esto acerca el mecanismo a MAR, lo que condiciona que
imputacion resulta defendible y debe declararse como supuesto.

\subsection{Valores centinela y fisicamente imposibles}\label{sec:dic:centinela}

Ademas de los faltantes explicitos, el consolidado contiene valores que no
pueden tomarse como mediciones: el codigo $-9999$ del \emph{datalogger} (26
celdas entre NOX, RAINF, RH, SR, TOUT, WSR y WDR), porcentajes de humedad
relativa mayores a 100\,\%, y lecturas fuera del rango fisico del sensor (por
ejemplo, TOUT $>50^{\circ}$C o WSR $>100$ km/h). Estos valores se convierten a
\texttt{NaN} antes de cualquier calculo, usando como filtro duro el
\textbf{rango de fabricante} de cada sensor (\cref{tab:dic-fabricante}), que es
fisico y estable. El rango de operacion anual que declara el SIMA no se usa
para este filtro porque es inconsistente entre anios (en 2022, por ejemplo,
declara para \texttt{PRS} un minimo mas alto que en 2021 y 2023, lo que
borraria datos validos); se trata en cambio como senial de sospecha a
revisar, no como regla de borrado automatico.

% -----------------------------------------------------------------------------
%  Tablas
% -----------------------------------------------------------------------------

\begin{table*}[t]
  \centering
  \caption{Contaminantes criterio. Todas son variables numericas continuas. El
    rango de operacion va del anio mas restrictivo al mas permisivo del
    periodo.}
  \label{tab:dic-contaminantes}
  \small
  \begin{tabular}{lllcccc}
    \toprule
    Variable & Descripcion & Unidad & Rango de operacion & Rango observado & Mediana & \% nulos \\
    \midrule
    \texttt{CO}    & Monoxido de carbono                   & ppm             & 0--8 a 0--20       & $0.0$--$37.0$    & 1.17  & 13.31 \\
    \texttt{NO}    & Monoxido de nitrogeno                 & ppb             & 0--350 a 0--500    & $0.3$--$945.1$   & 5.00  & 16.75 \\
    \texttt{NO2}   & Dioxido de nitrogeno                  & ppb             & 0--100 a 0--200    & $0.0$--$188.6$   & 11.40 & 17.37 \\
    \texttt{NOX}   & Oxidos de nitrogeno (NO $+$ NO$_2$)   & ppb             & 0--400 a 0--500    & $-9999$--$971.8$ & 17.30 & 16.82 \\
    \texttt{O3}    & \textbf{Ozono}: variable objetivo     & ppb             & 0--153 a 0--185    & $0.5$--$265.0$   & 23.00 & 14.73 \\
    \texttt{PM10}  & Particulas $<10\,\mu$m                & $\mu$g/m$^3$    & 0--800 a 0--999    & $1.0$--$1001.0$  & 49.00 & 4.63  \\
    \texttt{PM2.5} & Particulas $<2.5\,\mu$m               & $\mu$g/m$^3$    & 0--205.94 a 0--999 & $0.0$--$999.0$   & 17.00 & 25.23 \\
    \texttt{SO2}   & Dioxido de azufre                     & ppb             & 0--150 a 0--405    & $0.0$--$404.7$   & 3.80  & 15.03 \\
    \bottomrule
  \end{tabular}
\end{table*}

\begin{table*}[t]
  \centering
  \caption{Parametros meteorologicos. Todas son variables numericas continuas,
    con las salvedades de \texttt{WDR} (circular) y de \texttt{RAINF} y
    \texttt{SR} (cero inflado).}
  \label{tab:dic-meteo}
  \small
  \begin{tabular}{lllcccc}
    \toprule
    Variable & Descripcion & Unidad & Rango de operacion & Rango observado & Mediana & \% nulos \\
    \midrule
    \texttt{TOUT}  & Temperatura ambiente & $^{\circ}$C        & $-6.5$--45 a 0--45.5 & $-9999$--$785.0$ & 24.02  & 6.95  \\
    \texttt{RH}    & Humedad relativa     & \%                 & 0--100               & $-9999$--$725.0$ & 59.00  & 11.38 \\
    \texttt{SR}    & Radiacion solar      & kW/m$^2$           & 0--1 a 0--1.26       & $-9999$--$604.0$ & 0.00   & 3.66  \\
    \texttt{RAINF} & Precipitacion        & mm/h               & 0--25 a 0--80        & $-9999$--$360.0$ & 0.00   & 4.74  \\
    \texttt{PRS}   & Presion atmosferica  & mmHg               & 687.5--750           & $0.0$--$750.0$   & 714.40 & 5.08  \\
    \texttt{WSR}   & Velocidad del viento & km/h               & 0--35 a 0--75        & $-9999$--$337.0$ & 7.40   & 6.87  \\
    \texttt{WDR}   & Direccion del viento & grados azimutales  & 0--360               & $-9999$--$360.0$ & 116.00 & 8.35  \\
    \bottomrule
  \end{tabular}
\end{table*}

\begin{table}[t]
  \centering
  \caption{Catalogo de estaciones de monitoreo y su cobertura, calculada
    sobre el consolidado.}
  \label{tab:dic-estaciones}
  \small
  \begin{tabular}{llr}
    \toprule
    Clave & Cobertura & Horas \\
    \midrule
    \texttt{CE}   & 2020--2025                & 52\,605 \\
    \texttt{NE}   & 2020--2025                & 52\,595 \\
    \texttt{NE2}  & 2020--2025                & 52\,593 \\
    \texttt{NE3}  & \textbf{2021--2025}       & 43\,817 \\
    \texttt{NO}   & 2020--2025                & 52\,597 \\
    \texttt{NO2}  & 2020--2025                & 52\,594 \\
    \texttt{NO3}  & \textbf{dic.\ 2022--2025} & 27\,045 \\
    \texttt{NTE}  & 2020--2025                & 52\,595 \\
    \texttt{NTE2} & 2020--2025                & 52\,593 \\
    \texttt{SE}   & 2020--2025                & 52\,599 \\
    \texttt{SE2}  & 2020--2025                & 52\,595 \\
    \texttt{SE3}  & 2020--2025                & 52\,593 \\
    \texttt{SO}   & 2020--2025                & 52\,597 \\
    \texttt{SO2}  & 2020--2025                & 52\,593 \\
    \texttt{SUR}  & 2020--2025                & 52\,592 \\
    \bottomrule
  \end{tabular}
\end{table}

\begin{table}[t]
  \centering
  \caption{Porcentaje de valores faltantes por variable.}
  \label{tab:dic-faltantes}
  \small
  \begin{tabular}{lr@{\qquad}lr}
    \toprule
    Variable & \% nulo & Variable & \% nulo \\
    \midrule
    \texttt{PM2.5} & 25.23 & \texttt{RH}    & 11.38 \\
    \texttt{NO2}   & 17.37 & \texttt{WDR}   & 8.35  \\
    \texttt{NOX}   & 16.82 & \texttt{TOUT}  & 6.95  \\
    \texttt{NO}    & 16.75 & \texttt{WSR}   & 6.87  \\
    \texttt{SO2}   & 15.03 & \texttt{PRS}   & 5.08  \\
    \texttt{O3}    & 14.73 & \texttt{RAINF} & 4.74  \\
    \texttt{CO}    & 13.31 & \texttt{PM10}  & 4.63  \\
                   &       & \texttt{SR}    & 3.66  \\
    \bottomrule
  \end{tabular}
\end{table}

\begin{table}[t]
  \centering
  \caption{Rango de fabricante por parametro, adoptado como filtro duro de
    valores imposibles.}
  \label{tab:dic-fabricante}
  \small
  \begin{tabular}{lc}
    \toprule
    Parametro & Rango de fabricante \\
    \midrule
    PM10, PM2.5, O3      & $0$--$1000$     \\
    NO, NO2, NOx, SO2    & $0$--$500$      \\
    CO                   & $0$--$50$       \\
    RH                   & $0$--$100$      \\
    WS (\texttt{WSR})    & $0$--$180$      \\
    TEMP (\texttt{TOUT}) & $-50$--$50$     \\
    SR                   & $0$--$1.4$      \\
    BP (\texttt{PRS})    & $449.9$--$824.9$\\
    WD (\texttt{WDR})    & $0$--$360$      \\
    \bottomrule
  \end{tabular}
\end{table}


Este es un test